\documentclass[11pt]{article}
\usepackage[margin=1in]{geometry}
\usepackage{amsmath,amssymb,amsthm,mathtools}
\usepackage{bm}
\usepackage{hyperref}
\usepackage{enumitem}
\usepackage{booktabs}
\usepackage{array}
\usepackage{longtable}
\usepackage{graphicx}
\usepackage{setspace}
\usepackage{microtype}
\usepackage[T1]{fontenc}
\usepackage[utf8]{inputenc}

\title{Relational Actualism: Predictions, Departures, and Tests}
\author{Joshua Sandeman}
\date{\today}

\begin{document}
\maketitle

\begin{abstract}
Relational Actualism is not only an ontological and mathematical reconstruction of physics. It is also an exposed empirical program. This paper collects the current empirical consequences of the framework, separating direct discriminator predictions from broader structural departures and higher-order persistence consequences.
\end{abstract}

\section*{Front-matter status box}
\noindent\textbf{This paper does:} collect the framework's testable predictions; separate direct discriminator predictions from broader structural departures; and trace each major prediction back to the kernel, closure, implementation, and bridge layers on which it depends.

\section{Introduction}
Paper~I (\cite{RAKernelPaper}) defines the ontology, closure architecture, and implemented $d=4$ theory. Paper~II (\cite{RABridgePaper}) develops the strongest bridge results to known physics. The purpose of the present paper is different. It asks what follows empirically when one treats those results not as isolated achievements but as parts of a single process-based framework.

The paper is organized around four prediction classes: relatively robust, bridge-dependent, structural extension, and exploratory. A second organizing principle is that persistence is not the default state of things. Stable structure exists because graph growth is self-filtering: renewable motifs survive, cheap exits are blocked or narrowed, and increasingly organized systems persist only to the extent that their internal organization filters dominant disruption channels.

\section{Actualization-level predictions}
These predictions sit closest to the actualization program itself.

If actualization is fundamentally the irreversible inscription of a previously open possibility into the causal ledger, then there should exist a nontrivial threshold separating reversible possibility dynamics from record-forming actuality.

One of the most visible near-term discriminator classes concerns whether gravity alone can induce genuine quantum actualization. In the RA picture, if gravitational interaction remains below the relevant threshold for irreversible inscription, then entanglement-style signatures may occur without actualization.

\section{Gravity and cosmology departures}
The gravity bridge in Paper~II argues that implemented $d=4$ RA supports Einstein dynamics in the dense-regime continuum limit. This section concerns what happens when one moves away from that regime.

One characteristic structural consequence is that the framework does not begin with a primitive vacuum energy term. A second is that sparse-regime departures may mimic some apparent dark-sector anomalies.

\section{Matter and conservation predictions}
This section collects predictions tied to the particle-topology bridge, motif-renewal bridge, and the conservation structure of the implemented theory.

One of the strongest matter-sector consequences is exact baryon conservation as a topology-space persistence fact. The framework also motivates a comparatively austere matter sector and a structured strong-decay lifetime hierarchy tracking exit accessibility.

\section{Higher-order persistence: chemistry, life, and agency}
This section should be read more cautiously than the earlier ones. It extends a principle that has already appeared in the motif-lifetime program: stable structure persists when cheap disruption channels are filtered or blocked while renewal remains possible.

The causal-firewall idea is the most distinctive example. In disciplined form, it says that living systems may occupy regions of motif space whose internal organization blocks the cheapest thermal and chemical disruption pathways.

\section{Prediction dependency framework and summary}
Predictions are strongest when their dependency structure is visible. Every major prediction in the framework should, where possible, be traced through:
\begin{enumerate}[label=\arabic*.]
    \item kernel source;
    \item closure source;
    \item implementation source;
    \item bridge source;
    \item extension layer;
    \item failure mode.
\end{enumerate}

\begin{table}[htbp]
\centering
\caption{Compact summary of major prediction families in RA.}
\begin{tabular}{>{\raggedright\arraybackslash}p{3.1cm} >{\raggedright\arraybackslash}p{2.1cm} >{\raggedright\arraybackslash}p{2.2cm} >{\raggedright\arraybackslash}p{1.5cm} >{\raggedright\arraybackslash}p{4.2cm}}
\toprule
Prediction family & Main dependency & Prediction class & Status & What would count strongly against it \\
\midrule
Irreversibility threshold & actuality threshold + selectivity structure & RB & RB & no operational distinction between reversible possibility and irreversible inscription \\
BMV / gravity-actualization arena & actualization bridge to gravitational interaction & BR & BR & gravity-sensitive experiments show no threshold-sensitive difference at all \\
Direct strong-decay baseline & hadronic density closure + motif renewal & BR & BR & direct lifetime ordering fails under closure-backed density scale \\
Filtered lifetime hierarchy & $\sigma$-filters + pathwise exit + branching volume & BR/SE & BR & observed strong-decay hierarchy fails to track exit accessibility structure \\
No primitive $\Lambda$ term & ontology + gravity bridge & SE & SE & only viable cosmology requires fundamental vacuum term \\
Exact baryon conservation & topology closure + blocked exits & BR & BR & confirmed proton decay in a topology-forbidden channel \\
Dark-sector austerity / WIMP prohibition & finite motif universe + persistence logic & BR & BR & discovery of broad stable exotic particle sector outside plausible motif closure \\
Causal firewall / high-order persistence & recursive stabilization principle & SE/EX & EX & no empirical signature beyond metaphor, or no distinction between passive persistence and disruption filtering \\
\bottomrule
\end{tabular}
\end{table}

\section{Conclusion}
The strongest falsifiability lesson of the paper is this: RA should not be judged only by whether it can be made to sound deep. It should be judged by whether it survives contact with the specific failure modes it now names for itself.

If Relational Actualism is right, the universe is not a static inventory of things arranged in spacetime and governed by externally imposed laws. It is a self-writing growth process whose stable structures survive because graph dynamics filters their destruction.

\bibliographystyle{plain}
\bibliography{ra_suite_master_references}
\end{document}
